\documentclass[notheorems,aspectratio=169]{beamer}
\usetheme[mocha]{slidemint}
\usepackage{macromint}
\usepackage{figmint}

\slidemintsetup{
  footer-left = {\textcopyright~Author, 2026}
}

\title{Feature Check}
\subtitle{Title subtitle}
\author{Author}
\institute{Institute}
\date{2026}

\makeatletter
\newcommand*{\SlidemintPrintFontSize}[2]{%
  \begingroup
    #2%
    \typeout{SLIDEMINT_FIGMINT_FONT name=#1;size=\f@size}%
  \endgroup
}
\makeatother

\begin{document}

\typeout{SLIDEMINT_FIGMINT_THEME theme=\figminttheme}
\typeout{SLIDEMINT_ASPECT_RATIO width=\the\paperwidth;height=\the\paperheight}
\typeout{SLIDEMINT_FIGMINT_RATIO value=\figmintheighttowidthratio}
\SlidemintPrintFontSize{caption}{\usebeamerfont{caption}}
\SlidemintPrintFontSize{primary}{\figmintprimaryfont}
\SlidemintPrintFontSize{secondary}{\figmintsecondaryfont}

\maketitle

\section{Section Check}

\begin{frame}
  \frametitle{Text and Lists}
  \framesubtitle{Subtitle survives}

  \begin{itemize}
    \item \SlidemintTextbf{Strong text} and \SlidemintEmph{emphasis} share the theme colors.
    \item \SlidemintTextit{Italic text} keeps the selected font shape.
  \end{itemize}

  \begin{enumerate}
    \item First numbered item.
    \item Second numbered item.
  \end{enumerate}

  \begin{description}
    \item[GPU] Description label.
    \item[CPU] Description label with \href{https://example.org}{link text}.
  \end{description}

  \begin{ribbon}
    \begin{itemize}
      \item Ribbon item.
      \item \alert{Ribbon alert}.
    \end{itemize}
  \end{ribbon}
\end{frame}

\begin{frame}
  \frametitle{Blocks and Code}
  \framesubtitle{Theme environments}

  \begin{block}{Block}
    Normal block body.
  \end{block}

  \begin{alertblock}{Alert Block}
    Alert block body.
  \end{alertblock}

  \begin{exampleblock}{Example Block}
    Example block body.
  \end{exampleblock}
\end{frame}

\begin{frame}[fragile]
  \frametitle{Math and Code}
  \framesubtitle{Macromint and code colors}

  \[
    \set{x \in \nR^3 \suchthat \norm{x} = 1}
    \quad
    \gradient[\btheta]{\cL} + 123
  \]

  \begin{minted}{python}
def loss(theta):
    return norm(theta)
  \end{minted}

  \begin{minted}{text}
plain text code
  \end{minted}

  Inline code uses \mintinline{bash}|rsync| inside text.
\end{frame}

\begin{frame}
  \frametitle{Figures}
  \framesubtitle{Figmint follows the slide palette}

  \begin{figmintfigure}
    \begin{axis}[
      xmin=-0.5, xmax=2.5,
      ymin=0, ymax=1,
      xtick={0,1,2},
      nodes near coords,
    ]
      \addplot+[figmint bar] coordinates {(0,0.35) (1,0.65) (2,0.50)};
    \end{axis}
  \end{figmintfigure}
\end{frame}

\end{document}
